\section{Medición del tiempo}

\subsection{Wall clock vs tiempo de CPU}

En cualquier sistema operativo hay dos formas de medir cuánto tarda un programa:

\begin{itemize}
    \item \textbf{Wall clock}: el tiempo total transcurrido desde que el programa arranca hasta que termina. Es lo que un usuario percibe como ``cuánto tardó''. Incluye las esperas por E/S, las pausas por quantum y cualquier otra cosa que pase en el sistema.
    \item \textbf{Tiempo de CPU}: el tiempo que la CPU pasó efectivamente ejecutando nuestro programa. Se subdivide en tiempo de \textbf{usuario} (código nuestro) y tiempo de \textbf{sistema} (atendiendo llamadas al kernel).
\end{itemize}

Durante la materia, el profesor indicó que la métrica correcta para comparar rendimiento entre programas es el \textbf{wall clock}, porque refleja la experiencia real del usuario. El tiempo de CPU mide otra cosa (cuánto trabajó el procesador), pero no cuánto tardó el programa en completarse. En este trabajo se mide con wall clock en todos los casos.

\subsection{Qué opciones hay en C}

El lenguaje C ofrece varias funciones de tiempo. Esta tabla resume las que se consideraron:

\begin{center}
\scriptsize
\begin{tabular}{|l|l|c|}
\hline
\textbf{Función} & \textbf{Qué mide} & \textbf{Wall clock?} \\
\hline
\texttt{clock()} & Tiempo de CPU & \textcolor{red}{No} \\
\texttt{timespec\_get}\newline\texttt{(TIME\_UTC)} & Wall clock, pero afectado por cambios de hora & \textcolor{red}{No (frágil)} \\
\texttt{clock\_gettime}\newline\texttt{(CLOCK\_MONOTONIC)} & Wall clock monotónico & \textcolor{green}{Sí} \\
\texttt{clock\_gettime}\newline\texttt{(CLOCK\_MONOTONIC\_RAW)} & Wall clock, sin ajustes de NTP & \textcolor{green}{Sí, mejor} \\
\hline
\end{tabular}
\end{center}

\subsection{Por qué \texttt{CLOCK\_MONOTONIC\_RAW}}

De las cuatro, \texttt{clock\_gettime(CLOCK\_MONOTONIC\_RAW)} ofrece las mejores garantías para benchmarking:

\begin{itemize}
    \item Es \textbf{monotónica}: nunca devuelve un valor menor que el anterior, aunque se cambie la hora del sistema manualmente.
    \item Es \textbf{raw} (cruda): no se ve afectada por los ajustes graduales que NTP hace al reloj del sistema para sincronizarlo con servidores de hora externos. Estos ajustes pueden sumar o restar microsegundos al reloj, lo que introduce ruido en mediciones de corta duración.
\end{itemize}

\subsection{Qué se cronometra y qué no}

Se cronometra \textbf{solo el trabajo de multiplicación}. La reserva de memoria, la generación de las matrices aleatorias y la impresión de resultados quedan fuera del tiempo medido, porque no son el trabajo que estamos analizando. Esto tiene una consecuencia relevante para el perfilado: la función que contiene el triple bucle es la que domina el tiempo, y por eso es la que aparece en el reporte de GProf con el 100\% del tiempo propio.

\subsection{Por qué se repite el experimento 10 veces}

Una sola medición no es estadísticamente confiable: el sistema operativo puede tener otras tareas corriendo, una preempción puede caer justo en la corrida, o un proceso del sistema puede usar la CPU en el momento crítico. Para aislar este \textbf{ruido del sistema} se repite el experimento 10 veces por configuración y se calcula:

\begin{itemize}
    \item \textbf{Media} aritmética: el promedio de las 10 mediciones.
    \item \textbf{Mediana}: el valor central al ordenar las 10 mediciones. Es \textbf{robusta} frente a valores atípicos.
    \item \textbf{Desviación estándar} ($\sigma$): qué tanto se dispersan las mediciones.
\end{itemize}

En este trabajo se usa la \textbf{mediana} como estimador central del speedup, por la misma razón que en el caso de estudio 1.

\subsection{El orden de los bucles importa}

El script que ejecuta el experimento recorre las repeticiones por fuera y los tamaños por dentro:

\begin{verbatim}
for cada_T:
  for repeticion in 1..10:
    for cada_N:
      medir(N, T)
\end{verbatim}

El profesor explicó que si las 10 repeticiones de un mismo tamaño se hicieran seguidas, los datos quedarían \textit{calientes} en la caché del procesador y la medición saldría artificialmente rápida. Al intercalar tamaños entre repeticiones, cada medición se ve forzada a trabajar con datos nuevos en RAM, dando un resultado más realista.
